% Chapter 01: setup, self-similar reduction, main theorem.
% Source: gsqg_Nov26.tex, lines 185--529 (Introduction, Main Ideas, Notation) and
% the proof of thm:main (line ~664).

Throughout this chapter $\alpha\in(1,2)$ is fixed and we write $p=\alpha-1\in(0,1)$.
Unless said otherwise, all functions are real valued.  The Japanese bracket is
$\ang{y}=(1+y^2)^{1/2}$.  We set
\[
  c_\alpha=\frac{\Gamma(\alpha)\sin(\pi(\alpha-1)/2)}{\pi},
\]
which is a strictly positive real number for $\alpha\in(1,2)$ (since $\Gamma(\alpha)>0$ and
$0<\pi p/2<\pi/2$).  We write $W(y)=y_-=\max\{-y,0\}$ for the model corner.

%======================================================================
\section{Fourier conventions and the fractional multiplier}
%======================================================================

\begin{definition}[Model corner]
  \label{def:c01_corner}
  \lean{GSQG.corner}
  The model corner is $W:\R\to\R$, $W(y)=y_-:=\max\{-y,0\}$.
\end{definition}

\begin{definition}[Fourier transform]
  \label{def:fourier_convention}
  \lean{GSQG.fourier_convention, GSQG.fourierT, GSQG.invFourierT, GSQG.fourierMultiplier, GSQG.HasDistFourier, GSQG.distFourierInv, GSQG.IsPolyBounded}
  For $f\in L^1(\R;\C)$ the Fourier transform is
  \[
    \widehat f(\xi)=\int_{-\infty}^{\infty} f(x)e^{-ix\xi}\,dx ,\qquad \xi\in\R,
  \]
  extended to tempered distributions $\mathcal S'(\R)$ by duality,
  $\langle \widehat u,\varphi\rangle=\langle u,\widehat\varphi\rangle$.  The inverse transform is
  $\widecheck g(x)=\frac1{2\pi}\int_\R g(\xi)e^{ix\xi}\,d\xi$.  With this convention
  $\widehat{\p_x f}(\xi)=i\xi\widehat f(\xi)$.  A Fourier multiplier with symbol
  $m\in L^\infty_{\mathrm{loc}}(\R)$ of polynomial growth acts on $f\in\mathcal S(\R)$ by
  $m(D)f=(m\widehat f)^{\vee}$.  (Mathlib's \texttt{Real.fourierIntegral} uses
  $e^{-2\pi i x\xi}$; the two conventions differ by the substitution $\xi\mapsto 2\pi\xi$.)
\end{definition}

\begin{definition}[Secant]
  \label{def:c01_secant}
  \lean{GSQG.c01_secant, GSQG.c01_secantC}
  For $f:\R\to\R$ and $y\neq z$ put $S_f(y,z)=\dfrac{f(y)-f(z)}{y-z}$.  If $f\in C^1$ we extend
  by $S_f(y,y)=f'(y)$; then $S_f(y,z)=\int_0^1 f'(z+s(y-z))\,ds$ for all $y,z$.
\end{definition}

\begin{definition}[Admissible profile class]
  \label{def:c01_admissible_class}
  \lean{GSQG.c01_admissible_class}
  $\mathcal A$ denotes the set of $f\in C^1(\R)$ such that $f'\in L^\infty(\R)$ and $f'$ is
  locally Lipschitz.  Equivalently $f\in W^{2,\infty}_{\mathrm{loc}}(\R)$ with
  $f'\in L^\infty(\R)$ (taking the $C^1$ representative).  For $f\in\mathcal A$ and $R>0$
  write $\mathrm{Lip}_R(f')$ for the Lipschitz constant of $f'$ on $[-R,R]$.
  Note that $\mathcal A$ contains functions with linear growth, e.g.\ any $f$ with
  $f'(y)\to-1$ as $y\to-\infty$; no decay of $f$ is required.
\end{definition}

\begin{definition}[The fractional multiplier $|D|^{p}$ and $|D|^{p}\p_y$]
  \label{def:frac_multiplier}
  \lean{GSQG.frac_multiplier, GSQG.fracD, GSQG.fracDC, GSQG.fracDdy, GSQG.calpha}
  \uses{def:fourier_convention, def:c01_admissible_class}
  Let $0<p<1$, $\alpha=p+1$.
  \begin{enumerate}
    \item For $g\in L^\infty(\R)$ which is locally Lipschitz, define pointwise for every
      $y\in\R$
      \[
        |D|^{p}g(y):=c_\alpha\int_\R\frac{g(y)-g(z)}{|y-z|^{\alpha}}\,dz .
      \]
      (The integral is absolutely convergent by \cref{lem:c01_frac_well_defined}; no
      principal value is needed because $\alpha<2$.)
    \item For $f\in\mathcal A$ define $|D|^{p}\p_y f:=|D|^p(f')$, i.e.
      \[
        |D|^{\alpha-1}\p_y f(y)=c_\alpha\int_\R\frac{f'(y)-f'(z)}{|y-z|^{\alpha}}\,dz .
      \]
      This is the meaning of $|D|^{\alpha-1}\p_y f$ for non-decaying $f$ (e.g.\ $f'\to-1$ at
      $-\infty$): only $f'$ enters, and $f'$ is bounded.
    \item (Fourier side.) For $f\in\mathcal S(\R)$, $|D|^p f=(|\xi|^{p}\widehat f)^\vee$ and
      $|D|^{p}\p_y f=(i\xi|\xi|^{p}\widehat f)^\vee$; the two definitions agree on
      $\mathcal S(\R)$ by \cref{lem:c01_frac_singular_integral}.
  \end{enumerate}
\end{definition}

\begin{lemma}[Absolute convergence of the singular integral]
  \label{lem:c01_frac_well_defined}
  \lean{GSQG.c01_frac_well_defined}
  \uses{def:frac_multiplier}
  Let $0<p<1$, $\alpha=p+1$, and $g\in L^\infty(\R)$ be Lipschitz on $[y-1,y+1]$ with constant
  $\Lambda_y$.  Then $z\mapsto (g(y)-g(z))|y-z|^{-\alpha}$ is in $L^1(\R)$ and
  \[
    \bigl||D|^pg(y)\bigr|\le c_\alpha\Bigl(\frac{2\Lambda_y}{2-\alpha}
      +\frac{4\|g\|_{L^\infty}}{\alpha-1}\Bigr).
  \]
\end{lemma}
\begin{proof}
  Split $\R=\{|z-y|\le1\}\cup\{|z-y|>1\}$.  On the first set
  $|g(y)-g(z)|\,|y-z|^{-\alpha}\le\Lambda_y|y-z|^{1-\alpha}$, and
  $\int_{|h|\le1}|h|^{1-\alpha}dh=2/(2-\alpha)$ since $1-\alpha>-1$.  On the second set the
  integrand is bounded by $2\|g\|_\infty|y-z|^{-\alpha}$ and
  $\int_{|h|>1}|h|^{-\alpha}dh=2/(\alpha-1)$ since $\alpha>1$.
\end{proof}

\begin{lemma}[Affine invariance]
  \label{lem:c01_frac_affine}
  \lean{GSQG.c01_frac_affine}
  \uses{def:frac_multiplier}
  If $f_1,f_2\in\mathcal A$ and $f_1-f_2$ is affine, then $|D|^p\p_yf_1=|D|^p\p_yf_2$
  pointwise.  In particular $|D|^p\p_y f=|D|^p\p_y(f+y)$, so for $f$ with
  $f'(y)\to-1$ as $y\to-\infty$ one may compute $|D|^p\p_yf$ through the bounded function
  $f'+1$.
\end{lemma}
\begin{proof}
  $f_1'-f_2'$ is constant, so $f_1'(y)-f_1'(z)=f_2'(y)-f_2'(z)$ for all $y,z$.
\end{proof}

% MATHLIB-GAP
\begin{lemma}[Singular integral representation of $|\xi|^p$]
  \label{lem:c01_frac_singular_integral}
  \lean{GSQG.c01_frac_singular_integral}
  \uses{def:frac_multiplier, def:fourier_convention}
  Let $0<p<1$, $\alpha=p+1$.  For every $g\in\mathcal S(\R)$ and $y\in\R$,
  \[
    (|\xi|^p\widehat g)^\vee(y)=c_\alpha\int_\R\frac{g(y)-g(z)}{|y-z|^{\alpha}}\,dz,
    \qquad c_\alpha=\frac{\Gamma(1+p)\sin(\pi p/2)}{\pi}.
  \]
  Consequently, for $f\in\mathcal S(\R)$, $|D|^p\p_y f=(i\xi|\xi|^p\widehat f)^\vee$, and the
  linearization at $0$ of the gSQG nonlinearity is $|D|^{\alpha-1}\p_x$ (see
  \cref{lem:c01_linearization}).
\end{lemma}
\begin{proof}
  \uses{lem:c01_frac_well_defined}
  Write $z=y+h$ and symmetrize ($h\mapsto-h$, allowed since the integral converges absolutely
  by \cref{lem:c01_frac_well_defined}):
  \[
    c_\alpha\int_\R\frac{g(y)-g(y+h)}{|h|^{\alpha}}dh
    =\frac{c_\alpha}2\int_\R\frac{2g(y)-g(y+h)-g(y-h)}{|h|^{\alpha}}dh .
  \]
  Insert $g(x)=\frac1{2\pi}\int\widehat g(\xi)e^{ix\xi}d\xi$:
  $2g(y)-g(y+h)-g(y-h)=\frac1{2\pi}\int\widehat g(\xi)e^{iy\xi}\,2(1-\cos h\xi)\,d\xi$.
  Fubini is justified since $|\widehat g(\xi)|\,(1-\cos h\xi)|h|^{-\alpha}
  \le|\widehat g(\xi)|\min(h^2\xi^2,2)|h|^{-\alpha}$, which is integrable on $\R^2$
  ($\widehat g$ is Schwartz; $\int\min(h^2\xi^2,2)|h|^{-\alpha}dh=C|\xi|^{p}$).
  Substituting $s=|h\xi|$,
  \[
    \int_\R\frac{1-\cos h\xi}{|h|^{1+p}}dh=2|\xi|^p\int_0^\infty\frac{1-\cos s}{s^{1+p}}ds
    =2|\xi|^p\,\frac{\Gamma(1-p)\cos(\pi p/2)}{p}.
  \]
  The last integral is classical (integrate by parts once to get
  $\frac1p\int_0^\infty s^{-p}\sin s\,ds=\frac1p\Gamma(1-p)\sin(\pi(1-p)/2)$).
  Hence the right side equals $c_\alpha\frac{2\Gamma(1-p)\cos(\pi p/2)}{p}(|\xi|^p\widehat g)^\vee(y)$,
  and by $\Gamma(p)\Gamma(1-p)=\pi/\sin(\pi p)$ and $\sin(\pi p)=2\sin(\pi p/2)\cos(\pi p/2)$,
  \[
    \frac{p}{2\Gamma(1-p)\cos(\pi p/2)}=\frac{p\Gamma(p)\sin(\pi p)}{2\pi\cos(\pi p/2)}
    =\frac{\Gamma(1+p)\sin(\pi p/2)}{\pi}=c_\alpha .
  \]
  For the second claim apply the first to $g=f'$ and use $\widehat{f'}=i\xi\widehat f$.
\end{proof}

\begin{lemma}[Distributional meaning]
  \label{lem:c01_frac_duality}
  \lean{GSQG.c01_frac_duality}
  \uses{def:frac_multiplier, lem:c01_frac_singular_integral}
  Let $g\in L^\infty(\R)\cap W^{1,\infty}_{\mathrm{loc}}(\R)$ with $|g'(y)|\le C\ang{y}^N$ a.e.\ for
  some $C,N\ge0$.  Then $|D|^pg$ (pointwise, \cref{def:frac_multiplier}) is continuous, satisfies
  $\bigl||D|^pg(y)\bigr|\le C'\ang{y}^N$, and for every $\varphi\in\mathcal S(\R)$
  \[
    \int_\R |D|^pg(y)\,\varphi(y)\,dy=\int_\R g(y)\,|D|^p\varphi(y)\,dy ,
  \]
  where $|D|^p\varphi=(|\xi|^p\widehat\varphi)^\vee$ satisfies
  $|\,|D|^p\varphi(y)|\le C_\varphi\ang{y}^{-\alpha}$.  In particular the pointwise $|D|^pg$
  coincides with the distribution $\varphi\mapsto\langle g,|D|^p\varphi\rangle$, so a pointwise
  identity $\mc E_\epsilon(f)=0$ on $\R$ (\cref{def:residual}) for $f\in\mathcal A$ with $f''$ of
  polynomial growth implies $\mc E_\epsilon(f)=0$ in $\mathcal D'(\R)$.
\end{lemma}
\begin{proof}
  \uses{lem:c01_frac_singular_integral, lem:c01_frac_well_defined}
  The pointwise bound follows from \cref{lem:c01_frac_well_defined} with
  $\Lambda_y\le C\ang{|y|+1}^N\le C2^N\ang{y}^N$.  Continuity: dominated convergence in the
  representation $c_\alpha\int(g(y)-g(y+h))|h|^{-\alpha}dh$, with dominating function
  $\Lambda|h|^{1-\alpha}\mathbf 1_{|h|\le1}+2\|g\|_\infty|h|^{-\alpha}\mathbf 1_{|h|>1}$ uniform for
  $y$ in compacts.  The decay of $|D|^p\varphi$: by \cref{lem:c01_frac_singular_integral},
  $|D|^p\varphi(y)=c_\alpha\int(\varphi(y)-\varphi(z))|y-z|^{-\alpha}dz$; for $|y|\ge2$ split
  $|z-y|<|y|/2$ (use $|\varphi(y)-\varphi(z)|\le|y-z|\sup_{|w|\ge|y|/2}|\varphi'(w)|$, rapidly
  decaying) and $|z-y|\ge|y|/2$ (bounded by $\|\varphi\|_\infty\int_{|h|\ge|y|/2}|h|^{-\alpha}dh
  +(|y|/2)^{-\alpha}\|\varphi\|_{L^1}$).  Duality: for $\delta>0$ put
  $K_\delta(y,z)=c_\alpha|y-z|^{-\alpha}\mathbf 1_{|y-z|>\delta}$.  Then
  $\iint|g(y)-g(z)||\varphi(y)|K_\delta\,dz\,dy<\infty$ and likewise with $\varphi(z)$, so
  \[
    \iint (g(y)-g(z))\varphi(y)K_\delta
    =\tfrac12\iint(g(y)-g(z))(\varphi(y)-\varphi(z))K_\delta
    =\iint g(y)(\varphi(y)-\varphi(z))K_\delta
  \]
  (swap $y\leftrightarrow z$, symmetry of $K_\delta$).  Let $\delta\to0$: on the left the inner
  integral converges to $|D|^pg(y)$ and is dominated by $C'\ang{y}^N$, which is integrable
  against $|\varphi|$; on the right the inner integral converges to $|D|^p\varphi(y)$, dominated
  by $C_\varphi\ang{y}^{-\alpha}$, integrable against $|g|\le\|g\|_\infty$.  For the last
  statement apply this with $g=f'$; the other terms of $\mc E_\epsilon(f)$ are continuous
  functions.
\end{proof}

%======================================================================
\section{The gSQG front equation}
%======================================================================

\begin{definition}[The coefficient $A_\epsilon$]
  \label{def:A_eps}
  \lean{GSQG.A_eps, GSQG.A_eps_B}
  For $\epsilon\ge0$ and $u\in\R$ let $A_\epsilon(u)=(1+\epsilon u^2)^{-\alpha/2}$ and
  $B_\epsilon(u)=A_\epsilon(u)^{-1}=(1+\epsilon u^2)^{\alpha/2}$.  Then $A_\epsilon$ is smooth,
  $0<A_\epsilon\le1$, $A_\epsilon(0)=1$, $A_\epsilon'(0)=0$, $A_0\equiv1$, and
  $|A_\epsilon'(u)|=\alpha\epsilon|u|(1+\epsilon u^2)^{-\alpha/2-1}\le\alpha\epsilon|u|$.
\end{definition}

\begin{definition}[The nonlinearity $\mc J_\epsilon$]
  \label{def:c01_Jfull_eps}
  \lean{GSQG.c01_Jfull_eps}
  \uses{def:A_eps, def:c01_secant, def:c01_admissible_class}
  For $\epsilon\ge0$ and $f\in\mathcal A$ define, for every $y\in\R$,
  \[
    \mc J_\epsilon(f)(y)=c_\alpha\int_\R
    \frac{f'(y)-f'(z)}{\bigl((y-z)^2+\epsilon(f(y)-f(z))^2\bigr)^{\alpha/2}}\,dz
    =c_\alpha\int_\R\frac{f'(y)-f'(z)}{|y-z|^\alpha}A_\epsilon\bigl(S_f(y,z)\bigr)\,dz .
  \]
  (The two integrands agree for $z\ne y$ since
  $(y-z)^2+\epsilon(f(y)-f(z))^2=|y-z|^2(1+\epsilon S_f(y,z)^2)$.)  We write
  $\mc J:=\mc J_1$.  Note $\mc J_0(f)=|D|^p\p_yf$.
\end{definition}

\begin{lemma}[$\mc J_\epsilon$ is well defined]
  \label{lem:c01_J_well_defined}
  \lean{GSQG.c01_J_well_defined}
  \uses{def:c01_Jfull_eps, lem:c01_frac_well_defined}
  For $\epsilon\ge0$, $f\in\mathcal A$ and $y\in\R$, the integrand defining $\mc J_\epsilon(f)(y)$
  is absolutely integrable and
  $|\mc J_\epsilon(f)(y)|\le c_\alpha\bigl(\frac{2\Lambda_y}{2-\alpha}+\frac{4\|f'\|_\infty}{\alpha-1}\bigr)$,
  $\Lambda_y$ the Lipschitz constant of $f'$ on $[y-1,y+1]$.
\end{lemma}
\begin{proof}
  \uses{lem:c01_frac_well_defined}
  $0<A_\epsilon\le1$, so the integrand is bounded in absolute value by
  $|f'(y)-f'(z)||y-z|^{-\alpha}$; apply \cref{lem:c01_frac_well_defined} to $g=f'$.
\end{proof}

\begin{definition}[Front equation; classical solution]
  \label{def:gsqg_front_equation}
  \lean{GSQG.gsqg_front_equation}
  \uses{def:c01_Jfull_eps}
  Let $I\subset\R$ be an open interval.  A function $F:\R\times I\to\R$ is a
  \emph{classical solution of the gSQG front equation}
  \begin{equation}\label{eqn:front_equation}
    \p_tF(x,t)=c_\alpha\int_\R\frac{\p_xF(x,t)-\p_xF(x',t)}
    {\bigl((x-x')^2+(F(x,t)-F(x',t))^2\bigr)^{\alpha/2}}\,dx'
  \end{equation}
  on $\R\times I$ if
  \begin{enumerate}
    \item $F\in C^1(\R\times I)$, $F(\cdot,t)\in C^2(\R)$ for every $t\in I$, and $\p_x^2F$ is
      continuous on $\R\times I$;
    \item for every $t\in I$, $\p_xF(\cdot,t)\in L^\infty(\R)$ (hence $F(\cdot,t)\in\mathcal A$);
    \item for every $(x,t)\in\R\times I$, $\p_tF(x,t)=\mc J(F(\cdot,t))(x)$, the right side being
      the absolutely convergent integral of \cref{lem:c01_J_well_defined}.
  \end{enumerate}
  This is the graph (front) formulation of the gSQG patch problem with
  $\Omega(t)=\{x_2<F(x_1,t)\}$; we take \eqref{eqn:front_equation} as the definition and do not
  formalize its derivation from the gSQG system.
\end{definition}

\begin{lemma}[Agreement with the geometric front equation in the paper]
  \label{lem:c01_paper_front_equation}
  \lean{GSQG.c01_Jfull_paper_kernel, GSQG.c01_front_equation_paper}
  \uses{def:gsqg_front_equation, def:c01_Jfull_eps, def:c01_admissible_class}
  Let $0<p<1$, $\epsilon\ge0$, and $f\in\mathcal A$. For every $y\in\R$ the integral
  \[
    \mc J_\epsilon(f)(y)=c_\alpha\int_\R
      \frac{f'(y)-f'(z)}{((y-z)^2+\epsilon(f(y)-f(z))^2)^{(p+1)/2}}\,dz
  \]
  converges absolutely and equals the secant-factored operator used in Lean.
  Consequently every classical solution in \cref{def:gsqg_front_equation} satisfies the
  literal equation \eqref{eqn:front_equation}, with an absolutely convergent right-hand side.
\end{lemma}
\begin{proof}
  \uses{lem:c01_J_well_defined}
  For $y\ne z$ factor the denominator as
  $|y-z|^2(1+\epsilon S_f(y,z)^2)$; positivity permits the real-power product rule.
  The two integrands also agree at $y=z$ under Lean's total division convention, since the
  numerator vanishes. Transfer absolute convergence from \cref{lem:c01_J_well_defined}.
  For a classical solution, each spatial slice is $C^2$ with bounded first derivative,
  hence belongs to $\mathcal A$; apply the identity with $\epsilon=1$.
\end{proof}

\begin{lemma}[Linearization at the flat interface]
  \label{lem:c01_linearization}
  \lean{GSQG.c01_linearization}
  \uses{def:c01_Jfull_eps, def:frac_multiplier}
  Let $F\in\mathcal A$.  Then $\mc J_0(F)=|D|^{p}\p_xF$, and for every $x$,
  $\lim_{s\to0}s^{-1}\mc J(sF)(x)=|D|^p\p_xF(x)$.  More generally
  $\mc J(sF)=s\,\mc J_{s^2}(F)$ for $s\in\R$.
\end{lemma}
\begin{proof}
  \uses{def:c01_Jfull_eps, lem:c01_J_well_defined}
  Directly from \cref{def:c01_Jfull_eps}: $(sF)'=sF'$ and $S_{sF}=sS_F$, and
  $A_1(sv)=A_{s^2}(v)$.  The limit follows from dominated convergence, with the dominating
  function of \cref{lem:c01_J_well_defined}, since $A_{s^2}(S_F(x,z))\to1$ pointwise.
\end{proof}

\begin{lemma}[Symmetries of the front equation]
  \label{lem:c01_front_symmetries}
  \lean{GSQG.c01_front_symmetries}
  \uses{def:gsqg_front_equation}
  Let $F$ be a classical solution on $\R\times(a,b)$.  Then the following are classical solutions:
  \begin{enumerate}
    \item $F(x-x_0,t-t_0)+h_0$ on $\R\times(a+t_0,b+t_0)$, for $x_0,t_0,h_0\in\R$;
    \item $F_c(x,t)=c^{-1}F(cx,c^\alpha t)$ on $\R\times(c^{-\alpha}a,c^{-\alpha}b)$, for $c>0$;
    \item $\sq F(x,t)=F(-x,-t)$ on $\R\times(-b,-a)$.
  \end{enumerate}
\end{lemma}
\begin{proof}
  \uses{def:gsqg_front_equation}
  Regularity conditions (1)--(2) of \cref{def:gsqg_front_equation} are clearly preserved.
  (1): $\mc J$ only involves $F'$ and differences of $F$, and is translation covariant
  (substitute $x'\mapsto x'-x_0$).
  (2): $\p_tF_c=c^{\alpha-1}(\p_tF)(cx,c^\alpha t)$.  In $\mc J(F_c(\cdot,t))(x)$ we have
  $\p_xF_c(x)=(\p_xF)(cx)$, $(x-x')^2+(F_c(x)-F_c(x'))^2=c^{-2}\bigl((cx-cx')^2+(F(cx)-F(cx'))^2\bigr)$;
  substituting $w=cx'$, $dx'=dw/c$ gives $\mc J(F_c(\cdot,t))(x)=c^{\alpha-1}\mc J(F(\cdot,c^\alpha t))(cx)$.
  (3): $\p_t\sq F(x,t)=-(\p_tF)(-x,-t)$ and $\p_x\sq F(x,t)=-(\p_xF)(-x,-t)$; substituting
  $x'\mapsto -x'$ gives $\mc J(\sq F(\cdot,t))(x)=-\mc J(F(\cdot,-t))(-x)$.
\end{proof}

%======================================================================
\section{The self-similar ansatz and the profile equation}
%======================================================================

\begin{definition}[Self-similar ansatz]
  \label{def:self_similar_ansatz}
  \lean{GSQG.self_similar_ansatz, GSQG.self_similar_ansatz_L}
  For $T\in\R$ let $L_T(t)=(\alpha(T-t))^{1/\alpha}$ for $t<T$.  For $\epsilon>0$ and
  $f:\R\to\R$ define $F=F_{\epsilon,T,f}:\R\times(-\infty,T)\to\R$ by
  \begin{equation}\label{eqn:intro_self_similar_ansatz}
    F(x,t)=\sqrt\epsilon\,L_T(t)\,f\Bigl(\frac{x}{L_T(t)}\Bigr).
  \end{equation}
  We call $y=x/L_T(t)$ the self-similar variable and $f$ the profile.  The self-similar time is
  $\tau=-\log L_T(t)$.
\end{definition}

\begin{lemma}[Properties of $L$]
  \label{lem:c01_L_properties}
  \lean{GSQG.c01_L_properties}
  \uses{def:self_similar_ansatz}
  $L=L_T\in C^\infty((-\infty,T))$, $L>0$, $L(t)\to0$ as $t\uparrow T$, and $L'(t)=-L(t)^{1-\alpha}$.
\end{lemma}
\begin{proof}
  $L'=\frac1\alpha(\alpha(T-t))^{1/\alpha-1}\cdot(-\alpha)=-(\alpha(T-t))^{(1-\alpha)/\alpha}=-L^{1-\alpha}$.
\end{proof}

\begin{lemma}[Time derivative of the ansatz]
  \label{lem:c01_ansatz_time_derivative}
  \lean{GSQG.c01_ansatz_time_derivative}
  \uses{def:self_similar_ansatz, lem:c01_L_properties}
  Let $\epsilon>0$, $T\in\R$, $f\in C^1(\R)$ and $F=F_{\epsilon,T,f}$.  Then $F\in C^1(\R\times(-\infty,T))$ and, with $y=x/L(t)$,
  \[
    \p_xF(x,t)=\sqrt\epsilon f'(y),\qquad
    \p_tF(x,t)=-\sqrt\epsilon\,L(t)^{1-\alpha}\bigl(f(y)-yf'(y)\bigr).
  \]
  If $f\in C^k(\R)$ then $\p_x^jF(x,t)=\sqrt\epsilon L(t)^{1-j}f^{(j)}(y)$ for $j\le k$, and
  $F\in C^k(\R\times(-\infty,T))$.
\end{lemma}
\begin{proof}
  \uses{lem:c01_L_properties}
  Chain rule: $\p_t\bigl[L f(x/L)\bigr]=L'f(y)+Lf'(y)\cdot(-xL'/L^2)=L'(f(y)-yf'(y))$; use
  \cref{lem:c01_L_properties}.  Joint continuity of the derivatives follows from continuity of
  $f^{(j)}$, $L$, and $(x,t)\mapsto x/L(t)$.
\end{proof}

\begin{lemma}[Scaling of $\mc J$]
  \label{lem:c01_J_scaling}
  \lean{GSQG.c01_J_scaling}
  \uses{def:c01_Jfull_eps, lem:c01_J_well_defined}
  Let $\epsilon>0$, $\lambda>0$, $f\in\mathcal A$ and $G(x)=\sqrt\epsilon\,\lambda f(x/\lambda)$.
  Then $G\in\mathcal A$ and for every $x\in\R$
  \[
    \mc J(G)(x)=\sqrt\epsilon\,\lambda^{1-\alpha}\,\mc J_\epsilon(f)(x/\lambda).
  \]
\end{lemma}
\begin{proof}
  \uses{lem:c01_J_well_defined}
  $G'(x)=\sqrt\epsilon f'(x/\lambda)$, bounded and locally Lipschitz.  Put $y=x/\lambda$ and
  substitute $x'=\lambda z$, $dx'=\lambda dz$:
  $(x-x')^2+(G(x)-G(x'))^2=\lambda^2\bigl((y-z)^2+\epsilon(f(y)-f(z))^2\bigr)$, so the integrand
  becomes $\sqrt\epsilon(f'(y)-f'(z))\lambda^{-\alpha}(\dots)^{-\alpha/2}\cdot\lambda$.  The change of
  variables is legitimate for absolutely integrable functions (\cref{lem:c01_J_well_defined}).
\end{proof}

\begin{lemma}[Self-similar evolution]
  \label{lem:c01_self_similar_evolution}
  \lean{GSQG.c01_self_similar_evolution}
  \uses{def:gsqg_front_equation, def:self_similar_ansatz, lem:c01_J_scaling}
  Let $\epsilon>0$, $T\in\R$, and $f:\R\times\R\to\R$, $(y,\tau)\mapsto f(y,\tau)$, with
  $f,\p_yf,\p_y^2f,\p_\tau f$ continuous and $\p_yf(\cdot,\tau)$ bounded for each $\tau$.  Let
  $F(x,t)=\sqrt\epsilon L(t)f(x/L(t),-\log L(t))$, $L=L_T$.  Then $F$ is a classical solution on
  $\R\times(-\infty,T)$ iff for all $(y,\tau)$
  \begin{equation}\label{eqn:self_similar_front_equation_scaled}
    \p_\tau f=f-y\p_yf+\mc J_\epsilon(f(\cdot,\tau)).
  \end{equation}
\end{lemma}
\begin{proof}
  \uses{lem:c01_J_scaling, lem:c01_L_properties}
  With $\tau(t)=-\log L(t)$ we have $\tau'(t)=-L'/L=L^{-\alpha}$ (\cref{lem:c01_L_properties}).
  The chain rule gives
  $\p_tF=\sqrt\epsilon\bigl[L'(f-y\p_yf)+L^{1-\alpha}\p_\tau f\bigr]
  =\sqrt\epsilon L^{1-\alpha}\bigl[\p_\tau f-f+y\p_yf\bigr]$, and \cref{lem:c01_J_scaling} with
  $\lambda=L(t)$ gives $\mc J(F(\cdot,t))(x)=\sqrt\epsilon L^{1-\alpha}\mc J_\epsilon(f(\cdot,\tau))(y)$.
  Divide by $\sqrt\epsilon L^{1-\alpha}>0$.  Since $t\mapsto(\tau(t))$ is a bijection
  $(-\infty,T)\to\R$ and $(x,t)\mapsto(x/L,\tau)$ is a $C^\infty$ diffeomorphism of
  $\R\times(-\infty,T)$ onto $\R^2$, the identities are equivalent.  Regularity:
  $\p_xF=\sqrt\epsilon\p_yf$ and $\p_x^2F=\sqrt\epsilon L^{-1}\p_y^2f$ are continuous, $\p_xF(\cdot,t)$
  bounded.
\end{proof}

\begin{definition}[The lower-order nonlinearity $\sq{\mc J}_\epsilon$]
  \label{def:J_eps}
  \lean{GSQG.J_eps}
  \uses{def:A_eps, def:c01_secant, def:c01_admissible_class}
  For $\epsilon\ge0$, $f\in\mathcal A$ and $y\in\R$,
  \[
    \sq{\mc J}_\epsilon(f)(y)=c_\alpha\int_\R\frac{f'(y)-f'(z)}{|y-z|^\alpha}
    \Bigl(A_\epsilon\bigl(S_f(y,z)\bigr)-A_\epsilon\bigl(f'(y)\bigr)\Bigr)\,dz .
  \]
\end{definition}

\begin{definition}[The cubic remainder $\mc Q_\epsilon$]
  \label{def:c01_Q_eps}
  \lean{GSQG.c01_Q_eps}
  \uses{def:A_eps, def:c01_secant, def:c01_admissible_class}
  For $\epsilon\ge0$, $f\in\mathcal A$ and $y\in\R$,
  \[
    \mc Q_\epsilon(f)(y)=c_\alpha\int_\R\frac{f'(y)-f'(z)}{|y-z|^\alpha}
    \Bigl(A_\epsilon\bigl(S_f(y,z)\bigr)-1\Bigr)\,dz .
  \]
\end{definition}

\begin{lemma}[Decompositions of $\mc J_\epsilon$]
  \label{lem:c01_J_decomposition}
  \lean{GSQG.c01_J_decomposition}
  \uses{def:c01_Jfull_eps, def:J_eps, def:c01_Q_eps, def:frac_multiplier}
  For $\epsilon\ge0$ and $f\in\mathcal A$ the integrals defining $\sq{\mc J}_\epsilon(f)(y)$ and
  $\mc Q_\epsilon(f)(y)$ converge absolutely for every $y$, and
  \[
    \mc J_\epsilon(f)=|D|^{p}\p_yf+\mc Q_\epsilon(f)
    =A_\epsilon(f')\,|D|^p\p_yf+\sq{\mc J}_\epsilon(f)\qquad\text{pointwise on }\R .
  \]
\end{lemma}
\begin{proof}
  \uses{def:c01_Jfull_eps, lem:c01_frac_well_defined}
  Since $|A_\epsilon(S)-1|\le1$ and $|A_\epsilon(S)-A_\epsilon(f'(y))|\le1$, both integrands are
  dominated by $|f'(y)-f'(z)||y-z|^{-\alpha}\in L^1(dz)$ (\cref{lem:c01_frac_well_defined}).
  Write $A_\epsilon(S_f)=1+(A_\epsilon(S_f)-1)=A_\epsilon(f'(y))+(A_\epsilon(S_f)-A_\epsilon(f'(y)))$ in the
  second form of \cref{def:c01_Jfull_eps} and split the absolutely convergent integral by
  linearity; $A_\epsilon(f'(y))$ does not depend on $z$.
\end{proof}

\begin{lemma}[$\sq{\mc J}_\epsilon$ is of lower order]
  \label{lem:c01_Jtilde_vanishing}
  \lean{GSQG.c01_Jtilde_vanishing}
  \uses{def:J_eps, def:A_eps}
  Let $\epsilon\ge0$, $f\in\mathcal A$, $y\in\R$, and $\Lambda$ the Lipschitz constant of $f'$ on
  $[y-1,y+1]$.  Then for $|z-y|\le1$ the integrand of $\sq{\mc J}_\epsilon(f)(y)$ is bounded by
  $c_\alpha\,\alpha\epsilon\|f'\|_\infty\Lambda^2|y-z|^{2-\alpha}/2$.
\end{lemma}
\begin{proof}
  \uses{def:A_eps}
  $|f'(y)-f'(z)|\le\Lambda|y-z|$.  By the mean value theorem and \cref{def:A_eps},
  $|A_\epsilon(S_f(y,z))-A_\epsilon(f'(y))|\le\alpha\epsilon\|f'\|_\infty|S_f(y,z)-f'(y)|$ (note
  $|S_f|\le\|f'\|_\infty$), and $|S_f(y,z)-f'(y)|=|\int_0^1(f'(z+s(y-z))-f'(y))ds|\le\Lambda|y-z|/2$.
\end{proof}

\begin{definition}[The semilinear residual $\mc E_\epsilon$]
  \label{def:residual}
  \lean{GSQG.residual}
  \uses{def:frac_multiplier, def:A_eps, def:J_eps}
  For $\epsilon\ge0$ and $f\in\mathcal A$ define, pointwise for $y\in\R$,
  \[
    \mc E_\epsilon(f)(y)=|D|^{\alpha-1}\p_yf(y)+\frac1{A_\epsilon(f'(y))}
    \Bigl(f(y)-yf'(y)+\sq{\mc J}_\epsilon(f)(y)\Bigr).
  \]
\end{definition}

\begin{definition}[Quasilinear residual]
  \label{def:c01_quasilinear_residual}
  \lean{GSQG.c01_quasilinear_residual, GSQG.L0}
  \uses{def:frac_multiplier, def:c01_Q_eps}
  For $\epsilon\ge0$ and $f\in\mathcal A$ let
  $\sq{\mc E}_\epsilon(f)=|D|^{\alpha-1}\p_yf+f-yf'+\mc Q_\epsilon(f)$, i.e.\
  $\sq{\mc E}_\epsilon(f)=\mc L_0(f)+\mc Q_\epsilon(f)$ with $\mc L_0(f)=|D|^{\alpha-1}f'+f-yf'$.
\end{definition}

\begin{lemma}[Residual identities]
  \label{lem:c01_residual_identity}
  \lean{GSQG.c01_residual_identity}
  \uses{def:residual, def:c01_quasilinear_residual, lem:c01_J_decomposition}
  For $\epsilon\ge0$ and $f\in\mathcal A$, pointwise on $\R$:
  \[
    \sq{\mc E}_\epsilon(f)=f-yf'+\mc J_\epsilon(f)=A_\epsilon(f')\,\mc E_\epsilon(f).
  \]
  In particular, since $A_\epsilon>0$, for each $y$: $\mc E_\epsilon(f)(y)=0\iff
  f(y)-yf'(y)+\mc J_\epsilon(f)(y)=0$.
\end{lemma}
\begin{proof}
  \uses{def:residual, lem:c01_J_decomposition}
  First equality: \cref{lem:c01_J_decomposition}, first form.  Second: multiply
  \cref{def:residual} by $A_\epsilon(f'(y))$ and use the second form of \cref{lem:c01_J_decomposition}.
\end{proof}

\begin{definition}[Profile equation]
  \label{def:profile_equation}
  \lean{GSQG.profile_equation}
  \uses{def:residual}
  Let $\epsilon\ge0$.  A function $f\in\mathcal A$ solves the self-similar profile equation if
  \begin{equation}\label{eqn:main_ideas_self_similar_front_equation_semilinear}
    |D|^{\alpha-1}\p_yf+\frac{1}{A_\epsilon(\p_yf)}\Bigl(f-y\p_yf+\sq{\mc J}_\epsilon(f)\Bigr)=0
    \quad\text{i.e. }\ \mc E_\epsilon(f)(y)=0\ \ \forall y\in\R,
  \end{equation}
  together with the boundary conditions $\lim_{y\to-\infty}f'(y)=-1$ and
  $\lim_{y\to+\infty}f(y)=0$.  (The boundary conditions are not used in the reduction
  \cref{lem:self_similar_reduction}; they encode that $F$ approaches the corner $\sqrt\epsilon\,x_-$.)
\end{definition}

\begin{lemma}[Self-similar reduction]
  \label{lem:self_similar_reduction}
  \lean{GSQG.self_similar_reduction}
  \uses{def:gsqg_front_equation, def:self_similar_ansatz, def:residual}
  Let $\epsilon>0$, $T\in\R$ and $f\in C^2(\R)$ with $f'\in L^\infty(\R)$ (so $f\in\mathcal A$).
  Then $F=F_{\epsilon,T,f}$ is a classical solution of \eqref{eqn:front_equation} on
  $\R\times(-\infty,T)$ if and only if $\mc E_\epsilon(f)(y)=0$ for all $y\in\R$.
\end{lemma}
\begin{proof}
  \uses{lem:c01_ansatz_time_derivative, lem:c01_J_scaling, lem:c01_residual_identity, lem:c01_L_properties}
  Regularity: by \cref{lem:c01_ansatz_time_derivative} with $k=2$, $F\in C^2(\R\times(-\infty,T))$,
  in particular $F,\p_xF\in C^1$, and $\|\p_xF(\cdot,t)\|_\infty=\sqrt\epsilon\|f'\|_\infty<\infty$.
  Fix $t<T$, $x\in\R$, $y=x/L(t)$.  By \cref{lem:c01_ansatz_time_derivative},
  $\p_tF(x,t)=-\sqrt\epsilon L^{1-\alpha}(f(y)-yf'(y))$; by \cref{lem:c01_J_scaling} with
  $\lambda=L(t)$, $\mc J(F(\cdot,t))(x)=\sqrt\epsilon L^{1-\alpha}\mc J_\epsilon(f)(y)$.  Hence
  \[
    \p_tF(x,t)-\mc J(F(\cdot,t))(x)=-\sqrt\epsilon L(t)^{1-\alpha}\bigl(f(y)-yf'(y)+\mc J_\epsilon(f)(y)\bigr)
    =-\sqrt\epsilon L(t)^{1-\alpha}A_\epsilon(f'(y))\,\mc E_\epsilon(f)(y)
  \]
  by \cref{lem:c01_residual_identity}.  The prefactor is nonzero.  As $(x,t)$ ranges over
  $\R\times(-\infty,T)$ (with any fixed $t$), $y=x/L(t)$ ranges over all of $\R$, so the
  equation holds for all $(x,t)$ iff $\mc E_\epsilon(f)\equiv0$.
\end{proof}

%======================================================================
\section{The main theorem}
%======================================================================

\begin{theorem}[Self-similar corner formation]
  \label{thm:main}
  \lean{GSQG.main}
  \uses{def:gsqg_front_equation, def:self_similar_ansatz, def:profile_equation}
  Let $\alpha\in(\frac74,2)$, $p=\alpha-1$, $\beta=\frac{\alpha}{2(\alpha-1)}$.  There exist
  $\epsilon_0>0$ and $C>0$ (depending only on $\alpha$) such that for every
  $0<\epsilon<\epsilon_0$ there is a profile $f_\epsilon\in C^\infty(\R)$ with the following
  properties.
  \begin{enumerate}
    \item (Profile bounds.) For $y\le0$ and $0\le j\le2$:
      $|\p^j(f_\epsilon(y)+y)|\le C\ang{y}^{-p-j}$.  For $y\ge0$:
      $|f_\epsilon(y)|\le C\ang{y}^{-p}$ and $|\p^{j+1}f_\epsilon(y)|\le C\ang{y}^{-\beta+j/p}$ for
      $j\in\{0,1\}$.
    \item (Solution.) For every $T>0$, $F(x,t)=\sqrt\epsilon\,L(t)f_\epsilon(x/L(t))$,
      $L(t)=(\alpha(T-t))^{1/\alpha}$, is a classical solution of \eqref{eqn:front_equation}
      on $\R\times(-\infty,T)$ (\cref{def:gsqg_front_equation}), and $F\in C^\infty(\R\times(-\infty,T))$.
    \item (Corner convergence.) For all $T>0$ and $t<T$,
      \begin{equation}\label{eqn:intro_corner_convergence}
        \|F(\cdot,t)-\sqrt\epsilon\,x_-\|_{L^\infty(\R)}\le C\sqrt\epsilon\,(T-t)^{1/\alpha}.
      \end{equation}
    \item (Curvature blow-up.) There is $c=c(\epsilon)>0$ (independent of $T$ and $\delta$) such
      that for every $T>0$ and every $\delta>0$ there is $t_\delta<T$ with
      \begin{equation}\label{eqn:intro_curvature_blowup}
        \sup_{x\in(-\delta,\delta)}|\p_x^2F(x,t)|\ge c\,(T-t)^{-1/\alpha},\qquad t_\delta<t<T.
      \end{equation}
  \end{enumerate}
\end{theorem}
% NOTE: the paper writes the bounds with \lesssim; the constants are uniform in \epsilon because
% the constants in thm:fixed_point_approximation are uniform and ||r_eps|| <~ eps
% (thm:fixed_point_correction). In (4) we use sup over the open interval; since F(.,t) is C^2 this
% equals the L^\infty(-\delta,\delta) norm in the paper.
\begin{proof}
  \uses{lem:c01_profile_from_correction, lem:c01_profile_smooth, lem:c01_profile_bounds, lem:self_similar_reduction, lem:c01_ansatz_time_derivative, lem:c01_corner_convergence, lem:c01_curvature_blowup, lem:c01_profile_not_affine, thm:fixed_point_correction}
  Let $\epsilon_0$ be the minimum of the thresholds of \cref{lem:c01_profile_from_correction}
  (inherited from \cref{thm:fixed_point_correction}) and $1$, and for $0<\epsilon<\epsilon_0$ let
  $f_\epsilon=\ov f_{\epsilon,\mu_\epsilon}+r_\epsilon$ be the profile given there; it lies in
  $\mathcal A$ and $\mc E_\epsilon(f_\epsilon)\equiv0$.  By \cref{lem:c01_profile_smooth},
  $f_\epsilon\in C^\infty(\R)$.  (1) is \cref{lem:c01_profile_bounds}.  (2): by (1) $f_\epsilon'$ is
  bounded, so \cref{lem:self_similar_reduction} applies; smoothness of $F$ follows from
  \cref{lem:c01_ansatz_time_derivative} for all $k$.  (3) is \cref{lem:c01_corner_convergence} and (4) is
  \cref{lem:c01_curvature_blowup}, both applied with the bounds (1).
\end{proof}

\begin{lemma}[The profile solves the profile equation]
  \label{lem:c01_profile_from_correction}
  \lean{GSQG.c01_profile_from_correction}
  \uses{def:residual, def:parameters, def:approximate_solution_family, def:X_gamma}
  Fix $\alpha\in(\frac74,2)$ and the parameters $p,\beta,\gamma_1,\gamma_2$ of \cref{def:parameters}.
  There exists $\epsilon_1>0$ such that for every $0<\epsilon<\epsilon_1$, with $\mu_\epsilon$,
  $r_\epsilon\in X_{\gamma_2}$ given by \cref{thm:fixed_point_correction} and
  $\ov f_{\epsilon,\mu_\epsilon}$ by \cref{thm:fixed_point_approximation}, the function
  $f_\epsilon=\ov f_{\epsilon,\mu_\epsilon}+r_\epsilon$ belongs to $\mathcal A$, satisfies
  $\mc E_\epsilon(f_\epsilon)(y)=0$ for every $y\in\R$, and
  $|f_\epsilon''(y)|\le C\ang{y}^{1/p}$ for a.e.\ $y$, with $C$ independent of $\epsilon$.
  Moreover $\lim_{y\to-\infty}f_\epsilon'(y)=-1$ and $\lim_{y\to+\infty}f_\epsilon(y)=0$, so
  $f_\epsilon$ solves the profile equation (\cref{def:profile_equation}).
\end{lemma}
\begin{proof}
  \uses{thm:fixed_point_correction, thm:fixed_point_approximation, def:L_eps, def:N_eps}
  Take $\epsilon_1$ from \cref{thm:fixed_point_correction}.  Membership in $\mathcal A$:
  $\ov f=\ov f_{\epsilon,\mu_\epsilon}$ is $C^\infty$ with $|\ov f'|\le C$ on $\R$ (item 2 of
  \cref{thm:fixed_point_approximation}: $|\ov f'|\lesssim\ang{y}^{-\beta}$ for $y\ge0$ and
  $|\ov f'+1|\lesssim\ang{y}^{-p-1}$ for $y\le0$), and $r=r_\epsilon\in X_{\gamma_2}\subset
  W^{2,\infty}_{\mathrm{loc}}$ with $|r'|\le\|r\|_{X_{\gamma_2}}\ang{y}^{-\gamma_2+1/p}$,
  $|r''|\le\|r\|_{X_{\gamma_2}}\ang{y}^{-\gamma_2+2/p}$ a.e.; since $\gamma_2>\beta+1/p>1/p$
  these are bounded by $\|r\|_{X_{\gamma_2}}\ang{y}^{1/p}$, while $\ov f''$ has the stated
  polynomial growth and is bounded on every compact interval. So
  $f'$ is bounded and locally Lipschitz.
  The equation: by \cref{def:N_eps},
  $\mc N_{\epsilon,\ov f}(r)=\mc E_\epsilon(\ov f+r)-\mc E_\epsilon(\ov f)-\mc L_\epsilon(r)$ as functions on
  $\R$; inserting the correction equation $\mc L_\epsilon(r)=-\mc E_\epsilon(\ov f)-\mc N_{\epsilon,\ov f}(r)$
  gives $\mc E_\epsilon(\ov f+r)=0$ pointwise.  Limits: $|r(y)|\le\|r\|\ang{y}^{-\gamma_2}\to0$,
  $|r'(y)|\le\|r\|\ang{y}^{-\gamma_2-1}\to0$ as $y\to-\infty$, combined with item 1--2 of
  \cref{thm:fixed_point_approximation}.
\end{proof}

\begin{lemma}[Smoothness of the profile]
  \label{lem:c01_profile_smooth}
  \lean{GSQG.c01_profile_smooth}
  \uses{def:residual}
  The profile $f_\epsilon$ of \cref{lem:c01_profile_from_correction} is in $C^\infty(\R)$; hence
  also $r_\epsilon=f_\epsilon-\ov f_{\epsilon,\mu_\epsilon}\in C^\infty(\R)$ and the a.e.\ bounds on
  $r_\epsilon''$ coming from $X_{\gamma_2}$ hold everywhere.
\end{lemma}
\begin{proof}
  \uses{thm:fixed_point_smoothness, lem:c01_frac_duality, lem:c01_profile_from_correction}
  $f_\epsilon\in W^{2,\infty}_{\mathrm{loc}}$, $f_\epsilon'\in L^\infty$ and $|f_\epsilon''|\le C\ang{y}^{1/p}$,
  so by \cref{lem:c01_frac_duality} the pointwise identity $\mc E_\epsilon(f_\epsilon)\equiv0$ implies
  $\mc E_\epsilon(f_\epsilon)=0$ in $\mathcal D'(\R)$.  Apply \cref{thm:fixed_point_smoothness}.
  The bounds on $r_\epsilon''$ hold a.e.\ and $r_\epsilon''$ is continuous, hence they hold everywhere.
\end{proof}

\begin{lemma}[Profile bounds]
  \label{lem:c01_profile_bounds}
  \lean{GSQG.c01_profile_bounds}
  \uses{def:parameters}
  There exists $C>0$, independent of $\epsilon\in(0,\epsilon_1)$, such that the profile $f_\epsilon$
  of \cref{lem:c01_profile_from_correction} satisfies the bounds of item (1) of \cref{thm:main}.
\end{lemma}
\begin{proof}
  \uses{thm:fixed_point_approximation, thm:fixed_point_correction, def:X_gamma, lem:c01_profile_smooth}
  Let $\ov f=\ov f_{\epsilon,\mu_\epsilon}$, $r=r_\epsilon$, $\rho=\|r\|_{X_{\gamma_2}}\le C_0\epsilon\le C_0$
  (\cref{thm:fixed_point_correction}).  Recall $\beta+\frac1p<\gamma_2<2p+1$, so in particular
  $\gamma_2>p$ (because $\beta+1/p>1/p>1>p$) and $\gamma_2-\frac1p>\beta$.
  For $y\le0$ and $j\le2$: $|\p^j(\ov f+y)|\le C_j\ang{y}^{-p-j}$ by
  \cref{thm:fixed_point_approximation}, and $|\p^jr|\le\rho\ang{y}^{-\gamma_2-j}\le\rho\ang{y}^{-p-j}$.
  For $y\ge0$: $|\ov f|\le C\ang{y}^{-p}$ and $|r|\le\rho\ang{y}^{-\gamma_2}\le\rho\ang{y}^{-p}$;
  $|\ov f'|\le C\ang{y}^{-\beta}$ and $|r'|\le\rho\ang{y}^{-\gamma_2+1/p}\le\rho\ang{y}^{-\beta}$;
  $|\ov f''|\le C\ang{y}^{-\beta+1/p}$ and $|r''|\le\rho\ang{y}^{-\gamma_2+2/p}\le\rho\ang{y}^{-\beta+1/p}$.
  Add; all constants are uniform in $\epsilon$ and $\mu\in\ov B$ by
  \cref{thm:fixed_point_approximation}.  (Pointwise validity of the $r''$ bound:
  \cref{lem:c01_profile_smooth}.)
\end{proof}

\begin{lemma}[Bounded deviation from the corner]
  \label{lem:c01_profile_minus_corner}
  \lean{GSQG.c01_profile_minus_corner}
  Let $f:\R\to\R$ satisfy, for some $C>0$ and $p>0$, $|f(y)+y|\le C\ang{y}^{-p}$ for $y\le0$ and
  $|f(y)|\le C\ang{y}^{-p}$ for $y\ge0$.  Then $\sup_{y\in\R}|f(y)-W(y)|\le C$, where
  $W(y)=\max\{-y,0\}$.
\end{lemma}
\begin{proof}
  $W(y)=-y$ for $y\le0$ and $W(y)=0$ for $y\ge0$; and $\ang{y}^{-p}\le1$.
\end{proof}

\begin{lemma}[Corner convergence]
  \label{lem:c01_corner_convergence}
  \lean{GSQG.c01_corner_convergence}
  \uses{def:self_similar_ansatz}
  Let $\epsilon>0$, $T\in\R$, and $f:\R\to\R$ with $M:=\sup_y|f(y)-W(y)|<\infty$.  Then for all
  $t<T$,
  \[
    \sup_{x\in\R}\bigl|F_{\epsilon,T,f}(x,t)-\sqrt\epsilon\,x_-\bigr|=\sqrt\epsilon\,\alpha^{1/\alpha}M\,(T-t)^{1/\alpha}.
  \]
  In particular, under the bounds of \cref{thm:main}(1), \eqref{eqn:intro_corner_convergence} holds
  with constant $\alpha^{1/\alpha}C\le 2C$, and $F(\cdot,t)\to\sqrt\epsilon\,x_-$ uniformly as $t\uparrow T$.
\end{lemma}
\begin{proof}
  \uses{lem:c01_profile_minus_corner}
  $W$ is positively homogeneous of degree one: $\lambda W(y)=W(\lambda y)$ for $\lambda>0$.  Hence
  with $L=L_T(t)$ and $y=x/L$, $\sqrt\epsilon x_-=\sqrt\epsilon L\,W(y)$, so
  $F(x,t)-\sqrt\epsilon x_-=\sqrt\epsilon L\,(f(y)-W(y))$.  Since $x\mapsto x/L$ is a bijection of $\R$,
  the supremum over $x$ equals $\sqrt\epsilon L M$, and $L=\alpha^{1/\alpha}(T-t)^{1/\alpha}$.  Finally
  $\alpha^{1/\alpha}\le2$ for $\alpha\in(1,2)$, and use \cref{lem:c01_profile_minus_corner}.
\end{proof}

\begin{lemma}[The profile is not affine]
  \label{lem:c01_profile_not_affine}
  \lean{GSQG.c01_profile_not_affine}
  Let $f\in C^2(\R)$ with $\lim_{y\to-\infty}f'(y)=-1$ and $f$ bounded on $[0,\infty)$.  Then there
  exists $y_0\in\R$ with $f''(y_0)\ne0$.
\end{lemma}
\begin{proof}
  If $f''\equiv0$ then $f'$ is constant, and the limit forces $f'\equiv-1$, so $f(y)=b-y$ for
  some $b$, which is unbounded on $[0,\infty)$; contradiction.
\end{proof}
% NOTE: the hypotheses hold for f_eps by thm:main(1): |f'+1| <= C<y>^{-p-1} for y<=0 and
% |f| <= C<y>^{-p} for y >= 0. So the curvature blow-up does NOT need f''(0) != 0; but the
% constant c(eps) = sqrt(eps)|f''(y_0)| alpha^{-1/alpha} is non-explicit (y_0 obtained by
% contradiction) and no lower bound uniform in eps follows from the stated bounds.

\begin{lemma}[Curvature blow-up]
  \label{lem:c01_curvature_blowup}
  \lean{GSQG.c01_curvature_blowup}
  \uses{def:self_similar_ansatz}
  Let $\epsilon>0$ and $f\in C^2(\R)$ with $f''(y_0)\ne0$ for some $y_0\in\R$.  Put
  $c=\sqrt\epsilon\,\alpha^{-1/\alpha}|f''(y_0)|>0$.  Then for every $T\in\R$ and $\delta>0$, with
  \[
    t_\delta=\begin{cases}T-\frac1\alpha\bigl(\delta/|y_0|\bigr)^\alpha,& y_0\ne0,\\ T-1,& y_0=0,\end{cases}
  \]
  we have $\sup_{|x|<\delta}|\p_x^2F_{\epsilon,T,f}(x,t)|\ge c\,(T-t)^{-1/\alpha}$ for all
  $t_\delta<t<T$.  Moreover, if $f'$ is bounded, the geometric curvature
  $\kappa(x,t)=\p_x^2F/(1+(\p_xF)^2)^{3/2}$ satisfies
  $\sup_{|x|<\delta}|\kappa(x,t)|\ge c(1+\epsilon\|f'\|_\infty^2)^{-3/2}(T-t)^{-1/\alpha}$ for
  $t_\delta<t<T$.
\end{lemma}
\begin{proof}
  \uses{lem:c01_ansatz_time_derivative, lem:c01_L_properties}
  By \cref{lem:c01_ansatz_time_derivative}, $\p_x^2F(x,t)=\sqrt\epsilon L(t)^{-1}f''(x/L(t))$.  For
  $t_\delta<t<T$ we have $L(t)<\delta/|y_0|$ (if $y_0\ne0$), i.e.\ $x_0:=L(t)y_0$ satisfies
  $|x_0|<\delta$; if $y_0=0$ take $x_0=0$.  Then
  $|\p_x^2F(x_0,t)|=\sqrt\epsilon L(t)^{-1}|f''(y_0)|=c(T-t)^{-1/\alpha}$ since
  $L(t)^{-1}=\alpha^{-1/\alpha}(T-t)^{-1/\alpha}$.  For $\kappa$ use
  $|\p_xF|=\sqrt\epsilon|f'(x/L)|\le\sqrt\epsilon\|f'\|_\infty$.
\end{proof}

\begin{corollary}[Desingularization of a corner]
  \label{cor:c01_desingularization}
  \lean{GSQG.c01_desingularization}
  \uses{thm:main}
  Let $\alpha\in(\frac74,2)$, $0<\epsilon<\epsilon_0$, and $F$ the solution of \cref{thm:main} with
  $T=0$.  Then $\sq F(x,t)=F(-x,-t)$ is a classical solution of \eqref{eqn:front_equation} on
  $\R\times(0,\infty)$, $\sq F\in C^\infty$, and $\|\sq F(\cdot,t)-\sqrt\epsilon\,x_+\|_{L^\infty}\le
  C\sqrt\epsilon\,t^{1/\alpha}\to0$ as $t\downarrow0$, where $x_+=\max\{x,0\}$.
\end{corollary}
\begin{proof}
  \uses{lem:c01_front_symmetries}
  \Cref{lem:c01_front_symmetries}(3) with $(a,b)=(-\infty,0)$, and $(-x)_-=x_+$ in
  \eqref{eqn:intro_corner_convergence}.
\end{proof}

\begin{remark}[On the restriction $\alpha>7/4$]
  With $p=\alpha-1$, the parameter interval $\beta+\frac1p<\gamma<2p+1$ of \cref{def:parameters},
  i.e.\ $\frac{\alpha+2}{2(\alpha-1)}<\gamma<2\alpha-1$, is nonempty iff
  $p+3<4p^2+2p$, i.e.\ $(4p-3)(p+1)>0$, i.e.\ $\alpha>\frac74$.
\end{remark}
